\section{Proofs}
\begin{questions}
    % Direct Proof
    \item \begin{proof}
        Suppose $0 < x < 1$. Then $-1 < -x < 0$, meaning $0 < 1 - x < 1$. Therefore $x > 0$ and $1-x > 0$, so their product $x(1-x) > 0$.
    \end{proof}
    
    % Splitting into Cases
    \item \begin{proof}
        Suppose that $m$ and $n$ have the same parity. We split into two cases.
        \begin{itemize}
            \item If both $m$ and $n$ are even, then we may write $m = 2a$ and $n = 2b$ where $a$ and $b$ are integers. Hence
            \begin{align*}
                m + n &= (2a) + (2b) \\
                &= 2(a+b)
            \end{align*}
            which clearly means that $m + n$ is even.
            \item If both $m$ and $n$ are odd, then we may write $m = 2a + 1$ and $n = 2b + 1$ where $a$ and $b$ are integers. Hence
            \begin{align*}
                m + n &= (2a + 1) + (2b + 1)\\
                &= 2a + 2b + 2\\
                &= 2(a + b + 1)
            \end{align*}
            which clearly means that $m+n$ is even.
        \end{itemize}
        Hence, in both cases, $m + n$ is even.
    \end{proof}

    % Contrapositive Proof
    \item We consider a proof by contrapositive; the statement that we want to prove is ''if \textbf{not} ($a$ is even or $b$ is odd) then $a(b^2+5)$ is \textbf{not} even''. That is, ''if $a$ is \textbf{not} even \textbf{and} $b$ is \textbf{not} odd then $a(b^2+5)$ is \textbf{not} even'', meaning ''if $a$ is odd and $b$ is even then $a(b^2+5)$ is odd''.
    
    \begin{proof}
        Suppose that $a$ is odd and $b$ is even. Then we may write $a = 2m + 1$ and $b = 2n$ where $m$ and $n$ are integers. Hence
        \begin{align*}
            a(b^2+5) &= (2m+1)\left((2n)^2 + 5\right)\\
            &= (2m+1)(4n^2 + 5)\\
            &= 8mn^2 + 10m + 4n^2 + 5\\
            &= 8mn^2 + 10m + 4n^2 + 4 + 1\\
            &= 2(4mn^2 + 5m + 2n^2 + 2) + 1
        \end{align*}
        which clearly means that $a(b^2+5)$ is odd.
    \end{proof}
\end{questions}
